\label{sub:approachGMP}

The \textit{GNU Multi Precision Arithmetic Library (GMP)} allows implementing integral datatypes of arbitrary precision while supporting bit operations.
This means, one can choose the amount of bits used to represent a given value or directly assert a value of arbitrary size.
It is open source and scarcely documented, yet widely used due to its efficiency and general performance.
These values highlight the crucial role of GMP in this approach.

To use GMP in a project, one needs to install \code{libgmp\_dev} first.
The author used \code{apt} for this step.
Next, make sure to add the \code{-lgmp} flag during linking as it links the project with the GMP-library.
If the project is written in C++ also include GMP's C++ wrapper \code{-lgmpxx} as it provides overloaded operators frequently used in C++.
Since the author's code compiles using a \code{MAKEFILE}, the flag is set there.
Finally, include \code{gmp.h} in the files requiring GMP.
In the author's case it is solely included in \code{types.hpp}.

The following sections will use \code{mpz\_t} attributes frequently, which represent an integer of arbitrary size.
\code{mpz\_class} represent a class holding a single \code{mpz\_t} value providing the most common operators.
They, too, will be used rather frequently in the following subsections.

\subsubsection*{Approaches}\label{subsub:gmpApproaches}
When implementing this approach, two possible solutions were considered.
\begin{enumerate}
    \item Implement a wrapper struct for a single \code{mpz\_t} value.
    This struct can then be templated to match a certain bitsize.
    Assuming this struct is called \code{value\_wrapper<BIT\_SIZE>} this will result in \code{DPFnode} and \code{value\_t} being called \code{DPFnode = value\_wrapper<2*VALUE\_BITS>} and \code{value\_t = value\_wrapper<VALUE\_BITS>}.
    Then, however, one needs to take care of memory management manually, overload all operators for different datatypes and still is bound to the restrictions from \ref{sec:GeneralLimitations}.
    Also, drawbacks like communication issues between parties when using these structs from the second approach still exist.
    Therefore, this approach will not be pursued further\footnote{A basic implementation of this approach, however, can be found on the author's github at https://github.com/FelixWeik/PrivateRandomAccessComputations/tree/UnlimitedGMP}.
    \item Use the built-in \code{mpz\_class}.
    It takes care of memory management, initialization, and destruction automatically.
    One, however, cannot template its predefined bitsize.
    This means sizes will be arbitrarily chosen when, for example, a random value is chosen.
    Therefore, one can either choose if he limits the interval a number can be generated in or accept the arbitrary size of numbers.
    As this research is all about using values of arbitrary size, the author chose to go with the latter one.
    This makes this approach almost independent of macros like \code{VALUE\_BITS} (\ref{itemize:importantMacros}) as \code{\_\_m128i} vectors and sizes for communication-buffers can be generated at runtime.
    One can simply insert any value, the program will ultimately take care of it.
    Therefore, the author chose to further pursue and implement this approach\footnote{The implementation is located at https://github.com/FelixWeik/PrivateRandomAccessComputations/tree/gmp\_mpz\_class}.
\end{enumerate}

\subsubsection*{Applying GMP}\label{subsub:applyingGMP}
Again, main types, aliases, and structs are defined in \code{types.hpp} (\ref{subsub:implementationTypeshpp}).
\code{value\_t}, initially set as \code{uint64\_t}, is now replaced by a struct of the same name.
Implementing the datatype with a struct is necessary since GMP-values are set, initialized, and deleted during runtime and not like integral datatypes during precompilation.
Thus, one cannot directly use GMP-values as template arguments.
Since this property is required throughout the code the most common workaround uses a wrapper class or struct taking care of the initialization and deletion when called.
Then, one can use the struct for initializing templated structures in the code as required in, for example, \code{RDPFs}.

\code{value\_t} (\ref{itemize:importantAliases}), being a type alias, is now replaced by a struct.
This struct holds an attribute called \code{value} of type \code{mpz\_t}, which represents a signed integer.
As GMP offers no distinct type to implement unsigned integers, one simply adjusts the desired bitsize accordingly.
GMP adjusts the required bit-size by itself, so storing an equivalent of an \code{uint64\_t} still requires no more than 64-bits.
Since operations now take place on a struct instead of an integral datatype, one must overload the operators.
These operators call the already implemented GMP-operators on \code{value} and return the desired result.
Also, casting is necessary, for example when requiring an unsigned integer instead of \code{mpz\_t}.
Therefore, casting operators using GMP's built in casting functions are also required.
Finally, the constructors and destructors are implemented not to leak any memory.
In the following, the author will refer to structures of this type — which means structs that hold a single GMP-value and additionally implement overloaded operators, constructors, and destructors — as \textit{GMP-structs}.

Type-alias \code{nbits\_t} (\ref{itemize:importantAliases}), too, is used for initializing templated structures throughout the code.
Thus, in order to make it scalable, replacing it with a GMP-struct as above is necessary.
Recall, however, that \code{nbits\_t} solely counts the amount of bits required to store a value of type \code{value\_t}.
Leaving \code{nbits\_t} as is being \code{uint8\_t} would limit \code{value\_t} to a size of $2^8\text{ bits }\cong 32\text{ bytes}$.
Altering \code{nbits\_t} to equal \code{uint64\_t} will enable the user to store values reaching $2^{64} \text{ bit }\cong 2^{51}\text{ Terrabytes}$.
Of course, this value is far too big to achive testable and performant results.
This leads to the conclusion that scaling \code{nbits\_t} is of no great use in this research since it will solely impact performance and general memory usage.
Therefore, the author chose to set \code{nbits\_t} to \code{uint32\_t} in this approach which is the equivalent of storing a $512$ megabyte value.

Contrary to the prior approach, \code{VALUE\_BITS} stays untouched.
It defines the precision of the \code{mpz\_t} values.
No additional macros are required.

\paragraph{DPF-Nodes and Leaf-Nodes}\label{par:gmpDPFNodesAndLeafNodes}
\subsubsection*{Communication with GMP}\label{sub:CommunicationWithGMP}
\paragraph{Serialization and Deserialization}\label{par:gmpSerializationAndDeserialization}
